\chapter{Requirement Analysis}
\label{ch:requirements}

\section{Introduction}
This chapter converts the problem stated in Chapter~\ref{ch:problem} into a
specification. It describes how the requirements were elicited, identifies the
stakeholders and what each of them needs from the system, presents the use case
model, narrates the principal use cases in full, enumerates the functional and
non-functional requirements, records the hardware and software the system
depends on, and closes with a traceability matrix linking every requirement back
to the gap that motivated it and forward to the chapter in which it is realised.

\section{Requirement Elicitation Method}
The project had no external client, so requirements could not simply be
collected by interview. Four complementary techniques were used instead.

\textbf{Scenario construction.} Concrete situations were written down and used as
the primary source of functional requirements. Three of them shaped the interface
directly: a parent in Faisalabad deciding on Tuesday whether Thursday's outdoor
sports practice is safe; a construction supervisor in Multan deciding which of
the next three days to schedule outdoor work; and a person with asthma in Lahore
deciding whether to carry a mask tomorrow morning. Each scenario was traced
through the interface until it either produced an answer or exposed a missing
capability - the three-day aggregated view, the hazard banner and the ``cleanest
time of day'' question all originated this way.

\textbf{Analysis of the data.} A substantial exploratory study of the assembled
dataset, reported in Section~\ref{sec:design-eda}, was carried out before the
interface was designed. Several requirements came directly out of it. The finding
that January averages 146.8 against 91.5 in the cleanest month established that
the model would face strong recurring drift and therefore that drift monitoring
was necessary. The finding that the cities differ enormously - Faisalabad's mean
is roughly double Gilgit's - established that a city comparison view was
worthwhile and motivated the map.

\textbf{Adoption of published practice.} The operational requirements were taken
from the literature rather than invented. The production-readiness rubric of
Breck et al.~\cite{ref:mltestscore} supplied the requirements that the model be
compared against a simpler baseline, that training be reproducible, that model
quality be validated before serving and that rollback be possible. The
champion-challenger discipline~\cite{ref:siddiqi} supplied the promotion gate,
and the drift literature~\cite{ref:driftsurvey} supplied the monitoring
requirement.

\textbf{Iterative refinement from observed failure.} Because the lifecycle was
iterative, several requirements were discovered by running the system. The
requirement that ingestion be resumable and idempotent (FR-3) was not in the
original plan; it was added after a backfill froze partway through. The
requirement that the interface degrade gracefully when the backend is unavailable
(NFR-11) was added after the first deployment attempt served an interface with no
data behind it.

\section{Stakeholder Analysis}
Table~\ref{tab:stakeholders} identifies the stakeholders, their interest in the
system and the requirements that follow from that interest.

\begin{xltabular}{\textwidth}{P{2.9cm} P{4.4cm} Y}
\caption{Stakeholder analysis.}\label{tab:stakeholders}\\
\toprule
\bfseries Stakeholder & \bfseries Interest in the system & \bfseries Requirements arising \\
\midrule
\endfirsthead
\toprule
\bfseries Stakeholder & \bfseries Interest in the system & \bfseries Requirements arising \\
\midrule
\endhead
\bottomrule
\endlastfoot
General public (resident) & Wants to know whether the air will be safe over the
next few days, in their own city, in terms they can act on. & Multi-day forecast;
city selection; health category and advice rather than a bare number; a
prominent warning when the air will be dangerous; a plain-language explanation. \\
Vulnerable individuals & Children, the elderly and people with respiratory or
cardiac conditions face harm at lower exposure than the general population. &
Hourly resolution so that a safer time of day can be identified; explicit advice
per category; the earliest hour at which the unhealthy threshold is crossed. \\
Analyst or operator & Needs to know whether the system is still working and
whether its output can be trusted. & Published leaderboard and promotion history;
per-horizon accuracy; walk-forward backtest; drift status; realised forecast
error. \\
Assessor or reviewer & Needs to judge the engineering, not merely the output. &
Reproducibility; version history; documented architecture; unit tests; an audit
trail of model decisions. \\
Data provider (Open-Meteo) & Provides the observations free of charge and without
authentication. & Polite request behaviour with retry and backoff; no
unnecessary re-fetching; attribution in the interface. \\
Platform providers & Supply free scheduling, feature storage and hosting. &
Operation within free-tier limits; jobs that finish inside time caps; no
unbounded storage growth. \\
\end{xltabular}

\section{Use Case Diagram}
Figure~\ref{fig:usecase} presents the use case model. It distinguishes three
kinds of participant: the human actors who consume or operate the system, the
external data provider, and the scheduler, which is treated as a system actor
because it initiates the majority of the system's activity without any person
being present.

\fig{usecase}{0.92}{Use case diagram. Blue cases are public, green cases are
operational, and amber cases are automated and are initiated by the scheduler
rather than by a person.}{fig:usecase}

The most important structural feature of the diagram is that the four amber cases
have no human actor attached to them at all. This is the property that
distinguishes the system from a conventional application: its principal behaviour
- ingesting, retraining, gating, forecasting and monitoring - occurs on a
schedule, and the human actors interact with the results rather than with the
process.

\section{Use Case Narratives}

\subsection{View the Three-Day Forecast}
\begin{xltabular}{\textwidth}{Q{2.9cm} Y}
\toprule
\bfseries Use case ID & UC-1 \\
\bfseries Actor & Resident \\
\bfseries Goal & To learn what the air quality will be in a chosen city over the
next three days. \\
\bfseries Preconditions & A forecast document produced by a completed inference
run is available to the interface. \\
\bfseries Main flow & 1.~The user opens the dashboard. 2.~The system loads the
forecast document and displays the default city. 3.~The user selects a city from
the list of 22. 4.~The system displays the current index and its health
category, a chart of the next 72 hours with the 80\% prediction band, and the
category colour scale. 5.~The user switches to the daily view. 6.~The system
displays three calendar days, each with a mean index, a band and a category. \\
\bfseries Alternative flow & 3a.~The user instead selects a city by clicking its
marker on the map; the system updates the selection identically. \\
\bfseries Exception flow & 2a.~The forecast document cannot be retrieved. The
system displays an explanatory notice naming the failure rather than an empty
chart. \\
\bfseries Postcondition & The user has an hourly and a daily forecast, with
uncertainty, for the selected city. \\
\bfseries Requirements & FR-9, FR-10, FR-11, FR-19, NFR-3, NFR-11 \\
\bottomrule
\end{xltabular}

\subsection{Receive a Hazard Warning}
\begin{xltabular}{\textwidth}{Q{2.9cm} Y}
\toprule
\bfseries Use case ID & UC-2 \\
\bfseries Actor & Resident, in particular a vulnerable individual \\
\bfseries Goal & To be warned, in advance, that the air is forecast to become
dangerous, and to be told when. \\
\bfseries Preconditions & A forecast exists for the selected city. \\
\bfseries Main flow & 1.~During inference the system scans the city's 72-hour
curve. 2.~If the peak reaches the unhealthy threshold it classifies the episode
as a warning; if it reaches the very unhealthy threshold it classifies it as a
danger. 3.~It records the peak value, the hour of the peak, the category, the
standard advice for that category and the first hour at which the warning
threshold is crossed. 4.~The dashboard renders this as a banner above the
forecast whenever the severity is not none. \\
\bfseries Alternative flow & 2a.~The peak stays below the warning threshold; no
banner is shown and the forecast is presented normally. \\
\bfseries Postcondition & The user knows that an episode is coming, how bad it
will be, when it starts and what to do. \\
\bfseries Requirements & FR-12, FR-19, NFR-3 \\
\bottomrule
\end{xltabular}

\subsection{Understand Why the Forecast Is High}
\begin{xltabular}{\textwidth}{Q{2.9cm} Y}
\toprule
\bfseries Use case ID & UC-3 \\
\bfseries Actor & Resident \\
\bfseries Goal & To understand what is driving the forecast, rather than being
asked to accept a number. \\
\bfseries Preconditions & A champion model and a forecast for the city exist. \\
\bfseries Main flow & 1.~During inference the system identifies the hour with the
highest predicted index. 2.~It computes exact Shapley attributions for that
single prediction against the model's base value. 3.~It selects the five largest
contributions by absolute magnitude and translates each feature name into a
human-readable label. 4.~It composes a sentence stating the forecast value, the
base value and the signed contributions. 5.~The dashboard shows the sentence and
a signed bar for each contribution. \\
\bfseries Exception flow & 2a.~The attribution computation fails. The failure is
logged and the forecast is published without an explanation; the explanation card
is simply omitted. Explainability must never prevent a forecast from being
issued. \\
\bfseries Postcondition & The user can see which drivers pushed the forecast up
and which pushed it down, and by how much. \\
\bfseries Requirements & FR-13, FR-14, NFR-4, NFR-8 \\
\bottomrule
\end{xltabular}

\subsection{Run a What-If Scenario}
\begin{xltabular}{\textwidth}{Q{2.9cm} Y}
\toprule
\bfseries Use case ID & UC-4 \\
\bfseries Actor & Resident or analyst \\
\bfseries Goal & To interrogate the model by changing its inputs and observing
how the forecast responds. \\
\bfseries Preconditions & The live backend and the champion model are reachable. \\
\bfseries Main flow & 1.~The user opens the What-If tab. 2.~The system builds the
model input for the selected city at a lead time of 24 hours from live data and
returns the current real value of each of the five adjustable drivers, so that
the sliders begin at reality. 3.~The user adjusts one or more sliders.
4.~The user runs the simulation. 5.~The system predicts once with the unmodified
input and once with the overridden input and returns both values, both health
categories and the difference. 6.~The dashboard shows the two values side by
side, colour-coded by category, with the change. \\
\bfseries Exception flow & 2a.~The backend is unavailable, in which case the tab
is not offered at all rather than being offered and failing. \\
\bfseries Postcondition & The user has seen the model respond to a
counterfactual, clearly labelled as a model simulation rather than a causal
claim. \\
\bfseries Requirements & FR-15, FR-16, NFR-6, NFR-11 \\
\bottomrule
\end{xltabular}

\subsection{Retrain and Gate the Model}
\begin{xltabular}{\textwidth}{Q{2.9cm} Y}
\toprule
\bfseries Use case ID & UC-5 \\
\bfseries Actor & Scheduler (system actor) \\
\bfseries Goal & To incorporate new data into the model without allowing an
unattended run to degrade the service. \\
\bfseries Preconditions & The feature store holds the accumulated history; the
scheduled time has arrived. \\
\bfseries Main flow & 1.~The scheduler starts the training job. 2.~The job reads
the whole feature table. 3.~It constructs the multi-horizon supervised set and
subsamples it to the configured limit. 4.~It splits chronologically, holding out
the most recent fifth as validation. 5.~It scores the persistence baseline and
trains the three learned candidates. 6.~It selects the best fitted candidate by
validation RMSE as the challenger. 7.~It computes residual quantiles per horizon
for the prediction intervals. 8.~It compares the challenger with the reigning
champion. 9.~If the challenger is better it demotes the incumbent, records the
promotion and writes the bundle to the model registry; otherwise it records a
refusal and leaves the serving model untouched. 10.~It refreshes the evaluation
and monitoring reports and commits the artefacts. \\
\bfseries Alternative flow & 9a.~No champion exists yet, in which case the
challenger is promoted unconditionally as the first champion. \\
\bfseries Exception flow & 2a.~The feature store is unreachable; the job fails
loudly and the previous champion continues to serve, because promotion is the
only path by which the serving model can change. \\
\bfseries Postcondition & Either a strictly better model is in service, or the
previous model is still in service; in both cases the decision is recorded. \\
\bfseries Requirements & FR-4, FR-5, FR-6, FR-7, FR-8, NFR-1, NFR-7, NFR-9 \\
\bottomrule
\end{xltabular}

\subsection{Monitor Drift and Realised Error}
\begin{xltabular}{\textwidth}{Q{2.9cm} Y}
\toprule
\bfseries Use case ID & UC-6 \\
\bfseries Actor & Analyst; initiated by the scheduler \\
\bfseries Goal & To detect that the world has changed, and to know whether the
forecasts have actually been right. \\
\bfseries Preconditions & The feature store holds both historical and recent
rows; previous forecasts have been logged. \\
\bfseries Main flow & 1.~The system divides the feature history into a reference
window and the most recent fourteen days. 2.~For each monitored feature it bins
the reference distribution into deciles and computes the Population Stability
Index of the recent window against it. 3.~It classifies each feature as stable,
moderate or significant and identifies the worst. 4.~It loads the forecast log
and joins it to the observed index on city and target hour. 5.~It computes MAE
and RMSE over the joined rows and lists the five largest misses. 6.~It writes the
combined report, which the dashboard displays. \\
\bfseries Alternative flow & 4a.~No forecast has yet matured into an observation,
in which case the report states that scoring is waiting for actuals rather than
reporting a spurious zero. \\
\bfseries Postcondition & The operator can see both whether the inputs have moved
and whether accuracy has suffered. \\
\bfseries Requirements & FR-17, FR-18, NFR-9 \\
\bottomrule
\end{xltabular}

\subsection{Ask the Advisor a Question}
\begin{xltabular}{\textwidth}{Q{2.9cm} Y}
\toprule
\bfseries Use case ID & UC-7 \\
\bfseries Actor & Resident \\
\bfseries Goal & To ask a question in ordinary language and receive an answer
grounded in the system's real forecasts. \\
\bfseries Preconditions & The backend is reachable and an advisor credential is
configured. \\
\bfseries Main flow & 1.~The user types a question, or selects a suggested one.
2.~The backend sends the question to the language model together with the
declarations of four tools: list cities, get forecast, get history summary and
explain prediction. 3.~The model returns a structured tool call. 4.~The backend
executes the corresponding function against the real forecast data and returns
the result. 5.~Steps 3 and 4 repeat until the model produces a final answer, up
to a bounded number of rounds. 6.~The answer is returned and displayed. \\
\bfseries Exception flow & 1a.~No credential is configured; the endpoint returns
a clear, specific message and the interface reports that the advisor is
unavailable rather than failing silently. 5a.~The round limit is reached; the
system asks the user to rephrase. \\
\bfseries Postcondition & The user has an answer whose numbers came from the
system's own forecasts rather than from the language model's parameters. \\
\bfseries Requirements & FR-20, FR-21, NFR-5, NFR-11 \\
\bottomrule
\end{xltabular}

\section{Functional Requirements}
Table~\ref{tab:fr} enumerates the functional requirements. Priorities are
essential (E), important (I) or desirable (D).

\begin{xltabular}{\textwidth}{Q{1.5cm} Y Q{1.1cm}}
\caption{Functional requirements.}\label{tab:fr}\\
\toprule
\bfseries ID & \bfseries Requirement & \bfseries Pri. \\
\midrule
\endfirsthead
\toprule
\bfseries ID & \bfseries Requirement & \bfseries Pri. \\
\midrule
\endhead
\bottomrule
\endlastfoot
FR-1 & The system shall fetch hourly pollutant and weather observations for all
22 supported cities from a free public API, retrying transient failures with
exponential backoff. & E \\
FR-2 & The system shall engineer, per city and without cross-city contamination,
a feature table containing calendar, cyclical, wind-vector, lag, rolling and
change-rate features, and shall compute the EPA index as the target. & E \\
FR-3 & The system shall upsert engineered features into a feature store on a
composite key of city and timestamp, so that repeated or overlapping runs never
duplicate rows, and shall support a resumable historical backfill. & E \\
FR-4 & The system shall assemble a supervised dataset in which the forecast
horizon is an input feature and the target is the index observed at that horizon
ahead. & E \\
FR-5 & The system shall split training and validation data strictly
chronologically, never randomly. & E \\
FR-6 & The system shall train and score at least one baseline and three learned
candidate families in every training run, and shall report RMSE, MAE and $R^2$
for each. & E \\
FR-7 & The system shall promote a newly trained challenger to champion only if
it improves on the incumbent champion's validation RMSE, and shall record every
run, promoted or refused, in a persistent leaderboard. & E \\
FR-8 & The system shall store the promoted champion in a durable versioned model
registry that survives the destruction of the machine that trained it, and shall
load the champion from that registry at inference. & E \\
FR-9 & The system shall produce, for each supported city, a 72-hour hourly
forecast of the index and a three-day aggregation of it. & E \\
FR-10 & The system shall attach to every forecast point a prediction interval
derived from the empirical residual quantiles for that horizon. & E \\
FR-11 & The system shall label every forecast value with its EPA health category
and the standard advice for that category. & E \\
FR-12 & The system shall raise a graded hazard alert when a city's forecast peak
reaches the unhealthy or very unhealthy threshold, reporting the peak value, the
hour of the peak and the first crossing of the threshold. & E \\
FR-13 & The system shall compute Shapley attributions for the peak forecast hour
of each city and render the leading contributions as a plain-language sentence. & I \\
FR-14 & The system shall produce a global feature-importance ranking from
aggregated Shapley values. & D \\
FR-15 & The system shall expose the current real value of each adjustable driver
so that a what-if scenario begins from present conditions. & I \\
FR-16 & The system shall re-predict under user-supplied driver overrides and
return both the baseline and the scenario value with their categories. & I \\
FR-17 & The system shall compute the Population Stability Index per monitored
feature between a recent window and a historical reference, and classify each as
stable, moderate or significant. & I \\
FR-18 & The system shall log every issued forecast and later join it to the
observed index, reporting realised MAE, RMSE and the largest misses. & I \\
FR-19 & The system shall present the forecast, the map, the model leaderboard,
the per-horizon and backtest evaluation, and the monitoring report through a
public browser interface. & E \\
FR-20 & The system shall expose forecast, history and explanation functions as
declared tools so that a language model can answer questions from real data. & D \\
FR-21 & The system shall publish the same tools over the Model Context Protocol
for use by external clients. & D \\
FR-22 & The system shall serve an on-demand forecast for an arbitrary coordinate
pair, not only for the 22 pre-computed cities. & D \\
FR-23 & The system shall run ingestion hourly and retraining, inference,
evaluation and monitoring daily, without human initiation. & E \\
\end{xltabular}

\section{Non-Functional Requirements}
Table~\ref{tab:nfr} enumerates the non-functional requirements. Several of these
were as influential on the design as the functional requirements: NFR-1 and
NFR-2 between them dictated the entire architecture.

\begin{xltabular}{\textwidth}{Q{1.6cm} P{2.9cm} Y}
\caption{Non-functional requirements.}\label{tab:nfr}\\
\toprule
\bfseries ID & \bfseries Category & \bfseries Requirement \\
\midrule
\endfirsthead
\toprule
\bfseries ID & \bfseries Category & \bfseries Requirement \\
\midrule
\endhead
\bottomrule
\endlastfoot
NFR-1 & Cost & Every component shall operate within a free or negligible-cost
tier, so that the service can run indefinitely without a funding decision. \\
NFR-2 & Statelessness & No pipeline shall depend on state held on the machine
that ran a previous pipeline; all durable state shall live in the feature store,
the model registry or the repository. \\
NFR-3 & Performance & A page view shall not wait for a model to run: the served
forecast shall be a pre-computed document, read from disk. \\
NFR-4 & Explainability & No forecast shall be presented as a bare number; each
shall carry both an uncertainty band and, where computation permits, an
explanation. \\
NFR-5 & Groundedness & Any natural-language answer about air quality shall be
derived from the system's own forecast data through a tool call, never generated
from a language model's parametric knowledge. \\
NFR-6 & Honesty of framing & The what-if simulator shall be labelled explicitly
as a model simulation and shall not be presented as causal evidence. \\
NFR-7 & Reproducibility & Training shall be deterministic given the same data:
random seeds fixed, subsampling seeded, and the split defined by time rather than
by chance. \\
NFR-8 & Resilience & A failure in an auxiliary capability - explanation,
monitoring, forecast logging - shall be caught and logged and shall never
prevent the primary forecast from being produced. \\
NFR-9 & Auditability & Every model decision shall leave a durable record: the
leaderboard for promotions and refusals, the evaluation report for accuracy, the
monitoring report for drift and realised error. \\
NFR-10 & Portability & The storage layer shall be addressed through one facade
with interchangeable backends, so that the full pipeline can run locally on a
platform where the managed client cannot be installed. \\
NFR-11 & Graceful degradation & Where a capability requires a live backend, the
interface shall hide it rather than presenting it broken, and shall fall back to
committed data snapshots when no backend is configured. \\
NFR-12 & Maintainability & The codebase shall be organised into single-responsibility
modules with dependencies flowing in one direction, and the supported cities,
paths and horizons shall be defined in exactly one place. \\
NFR-13 & Usability & The interface shall communicate health status through the
standard EPA colour scale and category names, shall be readable on a phone, and
shall state when its data was generated. \\
NFR-14 & Politeness to the data source & Requests shall be retried with backoff
rather than repeated aggressively, historical windows shall be fetched in
chunks, and data already held shall not be re-fetched. \\
NFR-15 & Security & No credential shall appear in the repository; secrets shall
be supplied through environment variables locally and through the
continuous-integration secret store in automation. \\
\end{xltabular}

\section{Hardware and Software Requirements}
Table~\ref{tab:platform} records the platforms and libraries the system depends
on. The division between the development environment and the operating
environment is significant: the feature-store client cannot be installed on the
development platform at all, which is why the storage facade of NFR-10 exists.

\begin{xltabular}{\textwidth}{P{3.2cm} P{3.4cm} Y}
\caption{Hardware and software requirements.}\label{tab:platform}\\
\toprule
\bfseries Environment & \bfseries Component & \bfseries Requirement and purpose \\
\midrule
\endfirsthead
\toprule
\bfseries Environment & \bfseries Component & \bfseries Requirement and purpose \\
\midrule
\endhead
\bottomrule
\endlastfoot
Development & Workstation & A conventional laptop; no GPU is required, since the
production model is a gradient-boosted ensemble that trains on CPU in about a
minute. \\
Development & Python 3.11 & Pinned below 3.12 for dependency compatibility across
the machine-learning stack. \\
Development & Local Parquet store & Stands in for the feature store where the
managed client cannot be installed. \\
Automation & Linux runner & Ubuntu, Python 3.11, ephemeral; runs the hourly,
daily and manual workflows and is the only environment in which the feature-store
client is installed. \\
Automation & Scheduler & Cron-triggered workflows with concurrency groups, so
that two jobs never write to the feature store simultaneously. \\
Managed state & Feature store and model registry & Free serverless tier; holds
roughly 697,000 feature rows and the versioned champion bundles. \\
Serving & Container host & A small always-on container running the API from a
slim image; requires the OpenMP runtime for the gradient-boosting library. \\
Serving & Static edge host & Serves the compiled frontend bundle; no server
process. \\
Libraries (data) & pandas, requests, pyarrow & Tabular processing, HTTP with
retry, and columnar storage. \\
Libraries (models) & scikit-learn, XGBoost, SHAP, TensorFlow & Baselines and
metrics, the champion model, attributions, and the sequence-model study. \\
Libraries (serving) & FastAPI, Uvicorn, Pydantic & Typed asynchronous HTTP
interface with request validation. \\
Libraries (frontend) & React, Vite, Recharts, react-leaflet & Application
framework, build tooling, charting and cartography. \\
Libraries (testing) & pytest & Unit tests over feature engineering, the index
computation, the supervised builder and the promotion gate. \\
\end{xltabular}

\section{Requirements Traceability}
Table~\ref{tab:trace} traces each gap identified in Chapter~\ref{ch:problem}
through the requirements that address it to the chapter in which its realisation
is documented and the means by which it is verified.

\begin{xltabular}{\textwidth}{Q{1.0cm} P{3.3cm} P{2.6cm} Y}
\caption{Requirements traceability matrix.}\label{tab:trace}\\
\toprule
\bfseries Gap & \bfseries Requirements & \bfseries Realised in & \bfseries Verified by \\
\midrule
\endfirsthead
\toprule
\bfseries Gap & \bfseries Requirements & \bfseries Realised in & \bfseries Verified by \\
\midrule
\endhead
\bottomrule
\endlastfoot
G1 Nowcast only & FR-1, FR-4, FR-9, FR-23 & Sections
\ref{sec:impl-feature-pipeline} and \ref{sec:impl-inference} & Live forecast
documents containing 72 hourly points per city; endpoint checks. \\
G2 Narrow coverage & FR-1, FR-2, FR-4 & Sections \ref{sec:design-data} and
\ref{sec:design-mh} & 22 cities present in the feature store and in every
published forecast. \\
G3 Unreported accuracy & FR-5, FR-6, FR-7 & Section \ref{sec:impl-evaluate} &
Per-horizon table, per-horizon baseline comparison and five-fold walk-forward
backtest in Chapter~\ref{ch:testing}. \\
G4 No uncertainty & FR-10 & Section \ref{sec:design-intervals} & Interval bounds
present on every forecast point; band rendered in the interface. \\
G5 No explanation & FR-13, FR-14, FR-15, FR-16 & Sections \ref{sec:impl-shap} and
\ref{sec:impl-whatif} & Explanation object attached to each city's forecast;
what-if endpoint returning baseline and scenario. \\
G6 Silent decay & FR-3, FR-7, FR-8, FR-23 & Sections \ref{sec:impl-gate} and
\ref{sec:impl-cicd} & Leaderboard showing a promotion, a further promotion and a
refused tie; scheduled workflow history. \\
G7 No observation & FR-17, FR-18 & Section \ref{sec:impl-monitoring} & Published
PSI table with per-feature status; forecast log joined to observed values. \\
\end{xltabular}

\section{Summary of Requirement Analysis}
This chapter set out how the requirements were elicited - from scenarios, from
the data itself, from published operational practice and from observed failure -
and identified the stakeholders and what each needs. It presented the use case
model, whose defining feature is that the system's principal activity is
initiated by a scheduler rather than by a person, and narrated seven use cases in
full. It enumerated twenty-three functional and fifteen non-functional
requirements, recorded the platforms and libraries the system depends on, and
traced every gap from Chapter~\ref{ch:problem} through to its realisation and its
verification. The next chapter presents the design that satisfies this
specification.
